\chapter{Transformations et changements d'échelle}
\label{corriges:16}
\section*{Exercice 16.1 -- Images de points}
Les images sont respectivement $(2;-1)$, $(-2;-1)$, $(5;-1)$ et $(-1;2)$. Pour le quart de tour antihoraire, le déplacement de $2$ vers la droite devient $2$ vers le haut, et celui de $1$ vers le haut devient $1$ vers la gauche. Les quatre images sont à la même distance des autres points correspondants que dans la figure initiale.

\section*{Exercice 16.2 -- Rapport négatif}
Le facteur de longueur est $\abs{-1{,}5}=1{,}5$. Le périmètre devient $12$, l'aire $20\times1{,}5^2=45$ et le volume $40\times1{,}5^3=135$. Le signe négatif place les points images de l'autre côté du centre ; il ne rend aucune longueur, aire ou volume négatif.

\section*{Exercice 16.3 -- Retrouver une échelle}
Le facteur d'aire est $81/36=9/4$. Le facteur de longueur positif est sa racine : $r=3/2=1{,}5$. Un côté de $4\unite{cm}$ devient $6\unite{cm}$. Le facteur d'aire ne doit pas être appliqué directement aux côtés.

\section*{Exercice 16.4 -- Rosace et pavage}
L'angle est $360/8=45\degre$. Dans les triangles équilatéraux, chaque angle vaut $60\degre$ et $6\times60=360\degre$, ce qui complète le tour au sommet ; en les plaçant dans les secteurs successifs, les intérieurs ne se recouvrent pas. Un facteur $2$ multiplie les aires par $4$ : il ne les conserve pas.
